\section{Measurement, probability, and truth tables}

\subsection{Projective measurement}

For a measured qubit $q$ in a multi-qubit state,
\[
P(q=1)=\sum_{x:\,x_q=1}|a_x|^2,\qquad
P(q=0)=1-P(q=1).
\]
The simulator samples one outcome. Amplitudes incompatible with that outcome
become zero; compatible amplitudes are divided by the square root of the kept
probability so the state is normalized again. Measurement is destructive in
the sense that the pre-measurement superposition is generally unavailable
after collapse.

\subsection{Single runs and repeated runs}

One simulation run produces one measurement sample. To estimate a probability,
run the same preparation and circuit many times. If a plus state is measured
1 approximately 500 times in 1000 shots, the observed frequency $500/1000$
estimates the theoretical probability $1/2$. More shots reduce ordinary
sampling noise but never reveal a complex amplitude directly.

\subsection{Truth tables versus quantum amplitudes}

A classical truth table specifies output bits for listed input bits. It is a
good contract for arithmetic and reversible Boolean circuits operating on
computational-basis states. It cannot fully describe phase, interference,
entanglement, or a probability distribution.

The QPU truth-table runner prepares each input row, executes the circuit,
measures declared RETURNVALS, and compares those measured bits. In companion
\code{.qpuio} tables, cells are \code{0p}, \code{1p}, or \code{sp}. On an
expected output, \code{sp} means ``either measured bit is acceptable.'' Use
state equations and repeated experiments in addition to tables when a circuit's
purpose is genuinely quantum.

\subsection{A complete interference example}

Starting in zero, two Hadamards restore zero:
\[
\ket0\xrightarrow{H}\ket+
\xrightarrow{H}\ket0.
\]
Inserting Z changes the relative phase and changes the final answer:
\[
\ket0\xrightarrow{H}\ket+
\xrightarrow{Z}\ket-
\xrightarrow{H}\ket1.
\]
\begin{lstlisting}
MAIN-PROCESS Interference
CREATETOKEN -I Q
SET Q 0p
H -I Q -O Q
Z -I Q -O Q
H -I Q -O Q
MEASURE -I Q
RETURNVALS Q
\end{lstlisting}
The middle Z does not change an immediate 0/1 probability, yet interference at
the final H makes its effect deterministic. This is the central reason the
system documents matrices and phase as well as truth tables.

